%%%%%%%%%%  scan idx 354  |  folio 341  %%%%%%%%%%

Comme l'anneau considéré est intègre, on en déduit la relation suivante, valable par substitution dans tout \(\lambda\)-anneau, et pour tous éléments \(N,N'\) tels que \(\lambda^k(N)=0\) pour \(k>d\), \(\lambda^k(N')=0\) pour \(k>d'\)~;
\begin{equation}\tag{5.4.2}
  \lambda^p(N+N',x)\cdot\lambda_{-1}(N)=\lambda^p(N',x\cdot\lambda_{-1}(N)).
\end{equation}


\textbf{5.5.} Soit maintenant \(A\) un \(\lambda\)-anneau quelconque, et \(N\) un élément de \(A\) tel que \(\lambda^k(N)=0\) pour \(k>d\). On va associer à \(A\) un nouveau \(\lambda\)-anneau \(A_N\).

On définit d'abord une nouvelle multiplication sur \(A\) en posant pour \(x,y\in A\) :
\begin{equation}\tag{5.5.1}
  x\cdot_N y=x\cdot y\cdot\lambda_{-1}(N).
\end{equation}

L'anneau \(A\) muni de cette multiplication n'est plus unitaire; on considère donc l'anneau \(\mathbb Z\times A\), obtenu par adjonction d'une unité à \(A\). Ses opérations sont donc définies par :
\begin{equation}\tag{5.5.2}
\begin{aligned}
  (n,x)+(n',x')&=(n+n',x+x')\\
  (n,x)\cdot(n',x')&=(nn',nx'+n'x+x\cdot_N x').
\end{aligned}
\end{equation}
On considère \(A\) comme plongé dans cet anneau par \(x\rightsquigarrow(0,x)\). On définit alors une structure de pré-\(\lambda\)-anneau par :
\begin{equation}\tag{5.5.3}
  \lambda_t(n,x)=(1+t)^n\cdot\sum_{p=0}^{\infty}\lambda^p(N,x)\cdot t^p.
\end{equation}
C'est une structure de pré-\(\lambda\)-anneau en vertu de (5.4.1). On notera \(A_N\) le pré-\(\lambda\)-anneau ainsi obtenu. Pour voir que c'est un \(\lambda\)-anneau, il suffit par substitution de faire la démonstration dans le cas où \(A\) est le \(\lambda\)-anneau engendré sur \(\mathbb Z\) par les générateurs \(N,X,Y\), soumis aux relations \(\lambda^k(N)=0\) pour \(k>d\). On définit alors un homomorphisme \(A_N\to A\), en
